\documentclass[11pt]{article}

\usepackage[a4paper,margin=1in]{geometry}
\usepackage[T1]{fontenc}
\usepackage[utf8]{inputenc}
\usepackage{lmodern}
\usepackage{graphicx}
\usepackage{booktabs}
\usepackage{longtable}
\usepackage{hyperref}
\usepackage{xcolor}
\usepackage{listings}

\hypersetup{
  colorlinks=true,
  linkcolor=blue,
  urlcolor=blue,
  citecolor=blue
}

\lstset{
  basicstyle=\ttfamily\small,
  breaklines=true,
  columns=fullflexible,
  frame=single
}

\title{Chemistry-Aware Machine Translation Evaluation with Domain Terminology Verification}
\author{BASF AI}
\date{\today}

\begin{document}

\maketitle

\begin{abstract}
Chemical translation evaluation must distinguish general quality from retention of technical
terminology. We describe a segment-level methodology combining Google Patents Chemical abstract
pairs with an article-focused JRC-Acquis legal comparison corpus. The pipeline constructs source
snapshots, extracts reference-side terminology using complementary linguistic and learned approaches,
attaches external evidence, and selects exact reference spans through domain-specific large language
model curation. Translation generation separates a source-only baseline from an explicitly
reference-derived, oracle terminology condition. Evaluation combines corpus lexical overlap,
segment-level learned similarity, and strict and variant-aware term coverage. Intermediate artifacts
and provenance support controlled comparisons. Benchmark preparation is ongoing; comparative
translation results are not yet available. The protocol therefore makes no claim of improved MT
quality or chemically validated annotations.
\end{abstract}

\section{Introduction}
\label{sec:introduction}

Chemical translation requires preserving distinctions in names, materials, processes, formulas,
and quantities. General MT metrics assess overlap or learned similarity
\cite{papineni2002,post2018,rei2020,zhang2020}; these differ from checking whether a specific technical
term is retained correctly. We investigate how reference-side terminology annotations complement
conventional evaluation, and how providing such terminology affects a controlled translation comparison.

The primary domain is chemistry patent abstracts. JRC-Acquis legal articles offer a distinct,
terminology-sensitive comparison domain, not additional chemical training data. The evaluation unit
is a bounded segment, potentially containing several sentences, rather than a complete patent or
long legal document. The intended scale is a few hundred tokens; actual corpus lengths are reported
instead of assuming a shared 256-token mean.

The contributions are multilingual benchmark construction, provenance-preserving terminology
extraction and curation, and an evaluation protocol separating general quality from term retention.
These preserve the chemistry-aware evaluation objective. The infrastructure is implemented, but
the full translation comparison remains planned; no model ranking or superiority claim is made.

\section{Related Work and Resources}
\label{sec:resources}

JRC-Acquis comprises multilingual, predominantly legal EU texts \cite{steinberger2006}; OPUS provides
parallel distributions and tools \cite{tiedemann2012}. These aligned articles test transfer beyond
chemistry. Patent pairs come from a repository-documented Google Patents Chemical derivative
\cite{basfpatents}. Both supply evaluation text, unlike terminology databases.

BLEU measures clipped word $n$-gram overlap \cite{papineni2002}; SacreBLEU makes parameterization
explicit \cite{post2018}; chrF++ incorporates character and word overlap \cite{popovic2017}.
COMET estimates MT quality from source, hypothesis, and reference \cite{rei2020}, while BERTScore
aligns contextual hypothesis/reference embeddings \cite{zhang2020}. Explicit term coverage
complements these measures; semantic similarity is not chemical identity validation. Stanza supplies
linguistic analyses \cite{qi2020}. The configured XLM-R/NOBI extractor is English-trained
\cite{nobi}; its multilingual backbone does not establish equal extraction accuracy across languages.

\section{Data and Methodology}
\label{sec:methods}

\subsection{Evaluation Data and Preparation}
\label{sec:data}

\paragraph{Google Patents Chemical.}
The source is the BASF-AI within-document abstract-pairs dataset, documented
locally as within-publication abstract pairs from the CLIR Google Patents Chemical corpus
\cite{basfpatents}. The exporter normalizes whitespace, rejects empty text, retains accepted upstream
judge verdicts labelled \emph{exact} or \emph{high}, and selects in deterministic dataset order.
These are alignment-screening labels, not human chemical validation. The snapshot contains 2,750
pairs in 11 stored directions, capped at 250 per direction. Both-side approximate lengths of
128--384 tokens are preferred, followed by source-only, target-only, and out-of-range backfill;
the respective counts are 1,521, 626, 366, and 237. Length is a preference, not a hard criterion.

The active chemistry build retains German, English, Spanish, French, Russian, and Chinese,
omitting Japanese. Nine stored directions plus their source/target swaps yield 18 evaluation
directions and 4,500 planned segments. Reversal does not create independent publications.
The upstream split name \texttt{train} denotes the input partition, not training of our translation
systems. The export has no pinned upstream revision, so the local snapshot is needed for exact reuse.

\paragraph{JRC-Acquis articles.}
The article snapshot uses OPUS JRC-Acquis v3.0 aligned Moses files
\cite{steinberger2006,tiedemann2012}. It selects 250 document anchors shared across English,
Spanish, German, French, and Portuguese, generating all 20 ordered directions: 5,000 pairs,
250 per direction. Reverse rows are exact swaps; shared anchors do not imply identical passage
boundaries across every unordered pair. The separate definition snapshot is not included here.

After legacy markup cleanup and whitespace normalization, usable aligned units contain words,
3--180 approximate tokens on each side, and a maximum length ratio of 3. Units are concatenated
within document boundaries, targeting 450 source tokens with nominal bounds of 250--700 and at most
one chunk per document per pair. Strict filtering rejects residual markup, controls, all-capital
noise, and incomplete boundaries; it does not identify languages automatically. Recorded article
source lengths are 251--679, mean 418.6 approximate tokens, not 256. Counts use the repository
whitespace counter, not a provider tokenizer. This approximation is especially limited for Chinese;
the task remains segment-level despite this larger snapshot.

\paragraph{Partitioning and overlap.}
Both resources are evaluation inputs; no fine-tuning or finalized train/development/test partition
is specified. Small anchor pilots are pipeline checks, not independent held-out evidence. Before
final scoring, prompt-development and evaluation examples must be separated by JRC anchor and patent
publication, preferably family. Identifier and normalized-text duplicate audits, including reverse
directions and pilot overlap, remain pending. Public-corpus inclusion in proprietary pretraining
cannot be excluded; preserved identifiers facilitate audits but do not establish contamination-free data.

\subsection{Reference Terminology Annotation}
\label{sec:terminology}

Four automatic approaches operate on target references: domain-prompted LLM span proposals,
Stanza/Universal Dependencies syntactic spans, XLM-R/NOBI token-classified terms, and spaCy entities,
noun chunks, and filtered $n$-grams. Normalized duplicate keys merge provenance while retaining
surface forms; LLM proposals absent from the reference are rejected. The article configuration also
imports candidates from an earlier GLM-labelled manifest. This is reused automatic evidence, not
human annotation, and must be controlled in extractor ablations.

Chemistry attaches evidence from local IATE, Wikidata, and AGROVOC; articles use local IATE,
Wikidata, and UNTERM. These respectively supply EU terminology, multilingual labels/aliases,
FAO controlled vocabulary, and official UN terminology.
\footnote{\url{https://iate.europa.eu/};
\url{https://www.wikidata.org/wiki/Wikidata:Introduction};
\url{https://www.fao.org/agrovoc/about}; \url{https://unterm.un.org/unterm2/}.}
Scope and coverage differ: UNTERM's official languages exclude German and Portuguese. Local IATE
uses a CSV-derived SQLite index. Matches, variants, identifiers, and categories are stored as
evidence. Lookup misses or timeouts do not establish invalidity; chemistry network lookups have
a configured one-second timeout.

Domain-specific LLM curation selects numbered, enriched candidate IDs given the reference;
accepted terms remain exact reference spans. Both active configurations cap candidates at 80 and
refined terms at 16 per segment. They record \texttt{gpt-6-luna}, Responses API, temperature 0,
1,024 output tokens, disabled thinking, and reasoning effort \texttt{none}; this endpoint identifier
is not an independently verified immutable release. Chemistry uses four workers and a 180-second
LLM timeout; articles use two workers and 120 seconds. Manifests retain extractor provenance,
external support (\texttt{verified\_by}), and curation separately. Externally supported, LLM-selected,
and human-validated are distinct statuses; no expert gold terminology set or chemical-structure
validation is implemented. Chemical notation is preserved as text, not validated as molecular identity.

\subsection{Translation Generation}
\label{sec:generation}

The one-shot translator sends one segment, language labels, and a chemistry or legal prompt to
the selected provider, requesting only a translation and preservation of terms, formulas,
identifiers, quantities, units, and legal effect as appropriate. There is no fine-tuning,
retrieval augmentation, or constrained decoding. The planned paired comparison holds model and
decoding settings fixed while toggling manifest-term injection. A source-only baseline sees no
reference information. The other condition adds up to eight refined target terms with instructions
to use them only when the concept is present: soft hints, not enforced lexical constraints.

Since these hints derive from the evaluation reference, this is explicitly an \emph{oracle
terminology} condition, not deployable source-only glossary induction. Pilot configurations name
\texttt{gpt-4.1-nano} and \texttt{gpt-4.1-mini}; a legal run also names \texttt{gpt-6-luna}.
The final cross-domain roster and provider versions remain unspecified. Temperature 0 does not
guarantee deterministic hosted inference. Predictions, supplied terms, model metadata, and failures
are saved as JSONL. Dry-run source copying is a pipeline check, not an MT baseline.

\subsection{Quality and Terminology Evaluation}
\label{sec:evaluation}

The primary reporting set will combine corpus SacreBLEU and chrF++, segment COMET and BERTScore,
and strict/variant-aware target coverage. BLEU uses four word $n$-gram orders, case-sensitive
\texttt{13a} tokenization, and corpus scoring \cite{papineni2002,post2018}; chrF++ uses character
order 6, word order 2, and $\beta=2$ \cite{popovic2017}. They require hypotheses and references,
reward surface agreement, and can miss consequential chemical substitutions. Effective-order
sentence BLEU is diagnostic, not a substitute for corpus BLEU. This is not an official WMT25
scoring-protocol reproduction.

COMET uses source/hypothesis/reference with \texttt{Unbabel/wmt22-comet-da} \cite{rei2020};
BERTScore uses hypothesis/reference contextual token alignment with \texttt{xlm-roberta-large},
recording precision, recall, and F1 \cite{zhang2020}. These are computed per segment and averaged.
BERTScore precision/F1 concern embedding alignment, not chemical-term precision/F1; neither
metric certifies chemical identity, formula correctness, or legal effect.

For segment $i$, let $T_i$ contain applicable annotated target terms and $c_z(t)$ count occurrences
of $t$ in text $z$. Implemented strict coverage is
\begin{equation}
 C_i=\frac{100}{|T_i|}\sum_{t\in T_i}
 \min\left(\frac{c_{h_i}(t)}{c_{r_i}(t)},1\right),\quad c_{r_i}(t)>0,
\end{equation}
with hypothesis $h_i$ and reference $r_i$. Variant-aware coverage accepts explicit stored variants
in the numerator. Counting uses Unicode NFKC, case folding, whitespace normalization, and substring
matching, not word-boundary or chemical-entity matching. Case folding can erase meaningful chemical
case distinctions; extra hallucinated terms are not penalized. Report refined-term coverage and
its externally supported subset separately. Segments with no applicable terms are excluded, with
counts reported. Source-conditioned terminology success requires source-term alignment; target-only
annotations do not establish bilingual term accuracy.

Term-only BERTScore recall aligns the hypothesis against concatenated reference terms as a soft
diagnostic. The registry additionally supports reference-free COMETKiwi, XCOMET error spans, and a
domain-specific LLM MQM-style judge. The judge returns a 0--100 score and minor/major/critical
counts with weights 1/2/5; it is automatic, not expert annotation or an exact WMT FSP reproduction.
These remain supplementary pending access and calibration. BLEURT is optional and checkpoint-dependent;
NIST and terminology precision/F1 are not selected measures.

\subsection{Experimental Protocol and Reproducibility}
\label{sec:protocol}

Report domains and directions separately, macro-averaging direction means and recomputing corpus
overlap rather than averaging sentence BLEU. Paired comparisons use the same segments under baseline
and oracle prompts. Proposed ablations isolate extractor families, external evidence, and curation,
retaining exact-span validation. Annotation yield is not extraction accuracy without human labels;
cross-segment consistency and expert error analysis are not yet measured. Proposed paired resampling
groups by document/anchor and retains reversals together; statistical testing remains unimplemented.

Final scoring requires freezing the partition, model roster, metric subset, aggregation, and
uncertainty procedure. Reproduction artifacts include source checksums, Git revision, locked software,
extractor checkpoints/language models, exact prompts, provider identifiers and decoding settings,
dated IATE exports, and lookup caches. Selection, extraction, enrichment, curation, and manifest
checkpoints support stage-scoped hash reuse. Authentication/network failures are not completed
examples; successful-row denominators and failure counts must remain visible.

\section{Results and Analysis}
\label{sec:results}

Comparative translation results are not yet available. Planned analysis contrasts baseline and
oracle conditions by domain and direction, general quality with term coverage, and chemical-name,
formula, quantity, omission, and legal-effect errors. No empirical findings are claimed.

\section{Conclusion}
\label{sec:conclusion}

The methodology links bounded parallel segments to auditable reference terminology and complementary
evaluation measures. Its empirical value, oracle-term effects, and cross-domain generality remain
to be established by the planned comparisons.

\begin{thebibliography}{10}
\bibitem{papineni2002} K. Papineni, S. Roukos, T. Ward, and W.-J. Zhu.
Bleu: a Method for Automatic Evaluation of Machine Translation. ACL, 2002, pp. 311--318.
doi:10.3115/1073083.1073135.
\bibitem{post2018} M. Post. A Call for Clarity in Reporting BLEU Scores.
WMT, 2018, pp. 186--191. doi:10.18653/v1/W18-6319.
\bibitem{rei2020} R. Rei, C. Stewart, A. C. Farinha, and A. Lavie.
COMET: A Neural Framework for MT Evaluation. EMNLP, 2020, pp. 2685--2702.
doi:10.18653/v1/2020.emnlp-main.213.
\bibitem{zhang2020} T. Zhang, V. Kishore, F. Wu, K. Q. Weinberger, and Y. Artzi.
BERTScore: Evaluating Text Generation with BERT. 2020. arXiv:1904.09675.
\bibitem{steinberger2006} R. Steinberger, B. Pouliquen, A. Widiger, C. Ignat,
T. Erjavec, D. Tufis, and D. Varga. The JRC-Acquis: A multilingual aligned parallel corpus
with 20+ languages. LREC, 2006, pp. 2142--2147. arXiv:cs/0609058.
\bibitem{tiedemann2012} J. Tiedemann. Parallel Data, Tools and Interfaces in OPUS.
LREC, 2012, pp. 2214--2218. \url{https://aclanthology.org/L12-1246/}.
\bibitem{basfpatents} BASF-AI. AI4Chem CLIR Google Patents within-document abstract pairs.
Dataset artifact.
\href{https://huggingface.co/datasets/BASF-AI/ai4chem-clir-google-patents-within-document-abstract-pairs}{Hugging Face dataset record}.
\bibitem{popovic2017} M. Popovi\'{c}. chrF++: words helping character n-grams.
WMT, 2017, pp. 612--618. doi:10.18653/v1/W17-4770.
\bibitem{qi2020} P. Qi, Y. Zhang, Y. Zhang, J. Bolton, and C. D. Manning.
Stanza: A Python Natural Language Processing Toolkit for Many Human Languages.
ACL System Demonstrations, 2020, pp. 101--108. doi:10.18653/v1/2020.acl-demos.14.
\bibitem{nobi} tthhanh. XLMR Token Classifier for Term Extraction.
Model card, \texttt{xlm-ate-nobi-en-nes}.
\url{https://huggingface.co/tthhanh/xlm-ate-nobi-en-nes}.
\end{thebibliography}

\end{document}
